\section{Total positivity and Bessel-ratio tools used in Appendices~G--I}
\label{app:TP_tools}

This appendix collects the two analytic ``mechanisms'' used repeatedly in the one-plaquette A1 proofs
for $\Sp(2n)$, $\Spin(2n{+}1)$, and $\Spin(2n)$:

\begin{enumerate}[label=(\alph*)]
\item total-nonnegativity (TN) and a one-row shift inequality for Vandermonde-normalized minors;
\item explicit Amos-type bounds for adjacent modified-Bessel ratios.
\end{enumerate}

The purpose is (i) to avoid repeating the same TN and Bessel-ratio arguments in each family,
(ii) to make the logical dependencies in the A1 proofs auditably explicit.

\subsection{Modified Bessel functions: positivity, log-convexity, and Amos bounds}
\label{subsec:TP_tools_bessel}

We write $I_k(t)$ for the modified Bessel function of the first kind.
For integer $k\in\Z$ and $t>0$ one has $I_{-k}(t)=I_k(t)>0$.

\begin{lemma}[Bessel adjacent ratio monotonicity]
\label{lem:bessel_ratio_monotone}
Fix $t>0$ and define $q_k(t):=I_{k+1}(t)/I_k(t)$ for integers $k\ge 0$.
Then $k\mapsto q_k(t)$ is nondecreasing.
Equivalently, the sequence $k\mapsto I_k(t)$ is log-convex:
\begin{equation}
\label{eq:bessel_log_convex}
I_{k}(t)^2\ \le\ I_{k-1}(t)\,I_{k+1}(t)\qquad(k\ge 1).
\end{equation}
\end{lemma}

\begin{proof}
This is a classical property of $I_k(t)$ as a function of the order $k$ for fixed $t>0$.
One route is via total positivity: the bi-infinite Toeplitz matrix $(I_{i-j}(t))_{i,j\in\Z}$ is totally nonnegative
(Lemma~\ref{lem:PF_infty_bessel_toeplitz} below), and $2\times 2$ minors give \eqref{eq:bessel_log_convex}.
Monotonicity of $q_k(t)$ is equivalent to \eqref{eq:bessel_log_convex}.
\end{proof}

\begin{lemma}[Amos-type upper bound]
\label{lem:amos_upper_bound}
For every $t>0$ and every integer $k\ge 0$,
\begin{equation}
\label{eq:amos_upper_bound}
q_k(t)=\frac{I_{k+1}(t)}{I_k(t)}
\ \le\
\frac{t}{\,k+\tfrac12+\sqrt{t^2+(k+\tfrac12)^2}\,}.
\end{equation}
Equivalently,
\begin{equation}
\label{eq:q_as_exp_asinh}
q_k(t)\ \le\ \exp\Big(-\operatorname{arsinh}\big(\tfrac{k+\tfrac12}{t}\big)\Big).
\end{equation}
\end{lemma}

\begin{proof}
Inequality \eqref{eq:amos_upper_bound} is an explicit uniform upper envelope for Bessel ratios; see e.g.\ \cite{Hornik2014}.
The identity \eqref{eq:q_as_exp_asinh} follows from
$u+\sqrt{1+u^2}=e^{\operatorname{arsinh}(u)}$ applied with $u=(k+\tfrac12)/t$.
\end{proof}

\subsection{TN structure and a one-row shift inequality for normalized minors}
\label{subsec:TP_tools_TN_row_shift}

We use the following standard definition.

\begin{definition}[Totally nonnegative matrix]
\label{def:TN_matrix}
A real matrix $A$ (finite or infinite) whose rows and columns are indexed by totally ordered sets
is \emph{totally nonnegative (TN)} if every minor of every finite submatrix is nonnegative.
\end{definition}

\begin{lemma}[PF$_\infty$ for the Gross--Witten/Bessel coefficients]
\label{lem:PF_infty_bessel_toeplitz}
Fix $t>0$.
The bi-infinite Toeplitz matrix
\begin{equation}
\label{eq:bessel_toeplitz_matrix}
\mathcal T_t\ :=\ \big(I_{i-j}(t)\big)_{i,j\in\Z}
\end{equation}
is totally nonnegative.
\end{lemma}

\begin{proof}
We give a self-contained reduction to a finite-support Toeplitz criterion.

\medskip
\textbf{Step 1 (finite Laurent truncations).}
Fix $m\in\N$ and define the Laurent polynomial
\begin{equation}
\label{eq:GW_laurent_truncation_Fm}
F_m(z):=\Big(1+\frac{t}{2m}z\Big)^m\Big(1+\frac{t}{2m}z^{-1}\Big)^m
=\sum_{k=-m}^{m} a_k^{(m)} z^k.
\end{equation}
Let $\mathcal T_t^{(m)}:=(a_{i-j}^{(m)})_{i,j\in\Z}$ be the associated bi-infinite Toeplitz matrix.

\medskip
\textbf{Step 2 ($\mathcal T_t^{(m)}$ is TN for every $m$).}
The one-sided factor $P_m(z):=(1+\frac{t}{2m}z)^m$ has all its zeros equal to $-\frac{2m}{t}<0$.
Hence its coefficient sequence has a totally positive Toeplitz matrix by the finite-support Toeplitz criterion
\cite[Thm.~4.5]{Pinkus1996}.
The second factor is the reversed polynomial $z^{-m}P_m(z^{-1})$; its coefficient sequence also has a totally positive Toeplitz matrix.
Moreover, the Toeplitz matrix of the product $F_m=P_m\cdot z^{-m}P_m(z^{-1})$ is the product of the two Toeplitz matrices
(equivalently, the coefficients $a_k^{(m)}$ are the convolution of the two one-sided coefficient sequences),
and products of totally nonnegative matrices are totally nonnegative by Cauchy--Binet
(see e.g.\ \cite[Prop.~1.4]{Pinkus1996}).
Therefore $\mathcal T_t^{(m)}$ is totally nonnegative.

\medskip
\textbf{Step 3 (coefficientwise limit to $\mathcal T_t$).}
For $z=e^{i\theta}$ we have
\begin{equation}
\label{eq:GW_pointwise_limit_on_circle}
F_m(e^{i\theta})=\big|1+\tfrac{t}{2m}e^{i\theta}\big|^{2m}
=\Big(1+\frac{t}{m}\cos\theta+\frac{t^2}{4m^2}\Big)^m\xrightarrow[m\to\infty]{}e^{t\cos\theta}
\qquad(\theta\in[0,2\pi]).
\end{equation}
Also $0\le F_m(e^{i\theta})\le (1+\tfrac{t}{2m})^{2m}\le e^{t}$ for all $m$ and all $\theta$.
Since $F_m$ is a finite Laurent polynomial, its coefficients are its Fourier coefficients:
\begin{equation}
\label{eq:GW_coeff_as_Fourier}
a_k^{(m)}=\frac{1}{2\pi}\int_0^{2\pi} F_m(e^{i\theta})\,e^{-ik\theta}\,d\theta.
\end{equation}
Dominated convergence applied to \eqref{eq:GW_coeff_as_Fourier} and \eqref{eq:GW_pointwise_limit_on_circle} gives
$a_k^{(m)}\to \frac{1}{2\pi}\int_0^{2\pi} e^{t\cos\theta}e^{-ik\theta}d\theta$.
Using the standard cosine-Fourier representation
\(
I_k(t)=\frac{1}{\pi}\int_0^{\pi} e^{t\cos\theta}\cos(k\theta)\,d\theta
\) (which is used repeatedly in Appendices~G--I),
we obtain $a_k^{(m)}\to I_k(t)$ for every fixed $k\in\Z$.

Finally, any finite minor of $\mathcal T_t$ depends on finitely many entries $I_{i-j}(t)$.
For $m$ larger than the maximal index difference appearing in that minor, the corresponding minor of $\mathcal T_t^{(m)}$
uses only coefficients $a_k^{(m)}$ with $|k|\le m$ and converges entrywise to the desired minor of $\mathcal T_t$.
Since each $\mathcal T_t^{(m)}$ is TN, all such minors are nonnegative, and therefore the limit minor is nonnegative.
This proves that $\mathcal T_t$ is totally nonnegative.
\end{proof}

The next lemma is the TN ``engine'' behind all one-box inequalities in Appendices~G--I.

\begin{lemma}[One-row shift inequality for Vandermonde-normalized minors]
\label{lem:TN_row_shift}
Let $A$ be a TN matrix whose rows and columns are indexed by integers.
Fix an integer $N\ge 1$ and a strictly increasing $N$-tuple of column indices
\begin{equation}
\label{eq:TN_fixed_columns}
C=(c_1<c_2<\cdots<c_N).
\end{equation}
For every strictly increasing $N$-tuple of row indices
\begin{equation}
\label{eq:TN_rows_R}
R=(r_1<\cdots<r_N)
\end{equation}
define the Vandermonde factor
\begin{equation}
\label{eq:TN_vandermonde}
\Delta(R):=\prod_{1\le p<q\le N}(r_q-r_p)>0,
\end{equation}
and the normalized minor
\begin{equation}
\label{eq:TN_phi_def}
\Phi_A(R):=\frac{\det(A[R,C])}{\Delta(R)}\ \ge\ 0.
\end{equation}
Fix an index $i\in\{1,\dots,N\}$ such that replacing $r_i$ by $r_i+1$ and reordering preserves strict increase.
Let $R^+$ denote the resulting row tuple.
Then
\begin{equation}
\label{eq:TN_row_shift_bound}
\Phi_A(R^+)\ \le\ \frac{A[r_i+1,c_1]}{A[r_i,c_1]}\,\Phi_A(R).
\end{equation}
\end{lemma}

\begin{proof}
If $A[r_i,c_1]=0$ then $\det(A[R,C])=0$ (by TN and expansion along the first column), so \eqref{eq:TN_row_shift_bound} is trivial.
Assume henceforth that $A[r,c_1]>0$ for all rows $r$ appearing in $R$ and $R^+$.

Row-scaling by positive factors preserves total nonnegativity.
Let $D$ be the diagonal matrix with entries $D[r,r]=1/A[r,c_1]$ and set $\widetilde A:=DA$.
Then $\widetilde A$ is TN and satisfies $\widetilde A[\cdot,c_1]\equiv 1$ on the relevant rows.
To invoke the standard node-monotonicity statement from total positivity, we reduce to the strictly totally positive case:
by the Whitney density construction (see e.g.\ \cite[Thm.~2.6 and its proof]{Pinkus1996}),
every finite TN submatrix is the entrywise limit of strictly totally positive matrices.
Since the quantities in \eqref{eq:TN_row_shift_bound} depend continuously on the finitely many involved entries,
it suffices to prove the desired inequality for strictly totally positive matrices.

In the strictly totally positive case, for fixed columns $C$ the Vandermonde-normalized minors \eqref{eq:TN_phi_def} are generalized
(Popoviciu--Karlin) divided differences associated with the corresponding discrete totally positive (Chebyshev) system,
and they are monotone in each node.
For an explicit modern statement with numbered results and a proof via Sylvester's determinant identity, see
\cite[Cor.~3(ii)]{PalesRadacsi2018}; this monotonicity is traced there to \cite[Thm.~2]{Wasowicz2006}.
Applying this one-step node monotonicity to the single-node shift $r_i\mapsto r_i+1$ gives
\[
\Phi_{\widetilde A}(R^+)\le \Phi_{\widetilde A}(R).
\]
Undoing the row scaling uses that for any row tuple $R$ one has
\[
\det(A[R,C])=\Big(\prod_{p=1}^N A[r_p,c_1]\Big)\det(\widetilde A[R,C])
\quad\text{and hence}\quad
\Phi_A(R)=\Big(\prod_{p=1}^N A[r_p,c_1]\Big)\Phi_{\widetilde A}(R),
\]
and similarly for $R^+$. Therefore
\[
\begin{aligned}
\Phi_A(R^+)
&=\Big(\prod_{p=1}^N A[r_p^+,c_1]\Big)\Phi_{\widetilde A}(R^+)\\
&\le\Big(\prod_{p=1}^N A[r_p^+,c_1]\Big)\Phi_{\widetilde A}(R)\\
&=\frac{\prod_{p=1}^N A[r_p^+,c_1]}{\prod_{p=1}^N A[r_p,c_1]}
\,\Phi_A(R)\\
&=\frac{A[r_i+1,c_1]}{A[r_i,c_1]}\,\Phi_A(R).
\end{aligned}
\]
which is \eqref{eq:TN_row_shift_bound}.
\end{proof}

\subsection{TN of the classical-group Bessel moment matrices}
\label{subsec:TP_tools_moment_matrices}

Appendices~G--I use explicit moment matrices of the form
\begin{equation}
\label{eq:moment_matrix_template}
M_t(p,q)=\frac{2}{\pi}\int_0^{\pi} e^{t\cos\theta}\,\psi_p(\theta)\,\psi_q(\theta)\,d\theta,
\end{equation}
for specific trigonometric systems $\psi_p$ (sine with integer or half-integer frequencies, or cosine).

\begin{lemma}[TN of the moment matrices]
\label{lem:moment_matrix_TN}
Fix $t>0$.
Let $\psi_p(\theta)$ be one of the three families:
\begin{enumerate}[label=(\roman*)]
\item $\psi_p(\theta)=\sin(p\theta)$ for $p\in\N$ (symplectic case);
\item $\psi_p(\theta)=\sin\big((p+\tfrac12)\theta\big)$ for $p\in\N_0$ (odd orthogonal case);
\item $\psi_p(\theta)=\cos(p\theta)$ for $p\in\N_0$ (even orthogonal case).
\end{enumerate}
Then the infinite matrix $(M_t(p,q))$ defined by \eqref{eq:moment_matrix_template} is totally nonnegative.
\end{lemma}

\begin{proof}
The trigonometric kernels $(p,\theta)\mapsto \psi_p(\theta)$ in (i)--(iii) are strictly sign-regular kernels of all orders
on the ordered sets $p\in\N$ or $\N_0$ and $\theta\in(0,\pi)$.
A fundamental closure theorem in total positivity states that composing two sign-regular kernels against a strictly positive measure
preserves sign-regularity (hence TN at the discrete level).
Applying this to the composition
\(M_t(p,q)=\int \psi_p(\theta)\psi_q(\theta)\,e^{t\cos\theta}d\theta\)
with the strictly positive weight $e^{t\cos\theta}$ yields TN.
We refer to \cite[Vol.~I, Ch.~3]{Karlin1968} for the sign-regularity of the sine/cosine kernels and the composition theorem.
\end{proof}

\paragraph{How this appendix is used.}
In Appendices~G--I the matrices whose minors appear after Andr\'eief reduction are precisely of the form \eqref{eq:moment_matrix_template},
so Lemma~\ref{lem:moment_matrix_TN} supplies the TN hypothesis needed in Lemma~\ref{lem:TN_row_shift}.
Lemma~\ref{lem:amos_upper_bound} then provides the quantitative decay needed to turn the telescoping product over boxes
into the crossover estimate \eqref{eq:A1_canonical}.

